\subsection{Experimental setup details}
\label{app:setup-details}
\paragraph{Models, data, and calibration.}
The target is MiniCPM-o~4.5 (9B, based on Qwen3-8B). Matched controls
include raw backbones, a second MiniCPM generation, Qwen2.5-Omni,
Qwen2-Audio, and Freeze-Omni (Table~\ref{tab:models}). Representation
analyses use bf16, greedy decoding, and seed 42. The native benchmark
uses sampled decoding under the serving configuration
($T=0.7$, top-$k=20$); these are different experimental regimes.
The escalation expert is \texttt{gpt-5.5} with low reasoning effort.

The internal 600-query set contains SimpleQA traps
\citep{wei2024simpleqa}, MMLU-Pro knowledge questions
\citep{wang2024mmlupro}, TriviaQA facts \citep{joshi2017triviaqa},
Dolly/Alpaca chat prompts \citep{conover2023dolly,taori2023alpaca},
and GSM8K/MATH problems
\citep{cobbe2021gsm8k,hendrycks2021math,lightman2024lets}.
Its original 360/240 split and test composition remain fixed.
LLM judges score answer adequacy. A stronger re-judge of the
representation-study labels agrees on 96.2\% of text and 94.5\% of
audio cases; this agreement is not a validation of speech delivery.

Three calibration regimes must be distinguished: 360 examples for
single-position representation probes; 2,310 for the earlier
turn-based three-position gate; and 5,221 for the native gate, refitted
on native features and judged native-answer outcomes. The native fit
contains 4,979 core and 242 fresh training rows; thresholds use
five-fold out-of-fold scores from the core rows, grouped by language.
Conservative, balanced, and aggressive target nominal calibration
rates of 15/30/50\%; realized rates on new pools can differ.
Neither fresh-row scores nor external evaluation-pool quantiles enter
these thresholds. Seven rows overlapping external questions were
removed from the calibration set before refitting the gate and
recomputing its thresholds. All native-benchmark results reported here
come from reruns with this deduplicated fit
(Appendix~\ref{app:current-protocol}).

\paragraph{Native benchmark protocol and metrics.}
The MiniCPM upper block of Table~\ref{tab:transfer} evaluates the
internal test split and four external English speech-QA pools:
Speech TriviaQA, Speech Web Questions, and Speech Llama Questions
from OpenAudioBench
\citep{openaudiobench2025}, plus SD-QA recordings from VoiceBench
\citep{faisal2021sdqa,chen2024voicebench}. VoiceBench AlpacaEval is
reported separately on its 1--5 scale. The archived Mandarin Speech
Reasoning QA pool is retained in the appendix diagnostics and excluded
from the main-table macro average. External audio comes from the
benchmark releases; internal questions use a fixed TTS rendering.
Each query has a separate native-protocol session for each arm:
local-only, three gate tiers, and always-escalate. There is one
recorded run per query and arm, without repeated-run averaging.

Audio enters the local model in one-second chunks. At an escalation,
transcription and the expert produce a text answer, which is cleaned
and capped at 400 characters for teacher-forced TTS relay. We report
\emph{answer-content accuracy}: the judge sees local generated text
or the intended relay text, not a transcription of played audio.
The benchmark disables local audio generation. Uplink transcription
uses the final 30 seconds of the full prerecorded question: it can
access unconsumed audio at early onset and loses the beginning of
longer questions. Thus the experiment evaluates routing under this
input protocol, without establishing strictly causal interaction.

For the MiniCPM block, the reported random reference is
\[
 A_{\rm mix}(r)=(1-r)A_{\rm local}+rA_{\rm always},
\]
where $r$ is the gate's realized escalation rate. This follows
Table~\ref{tab:transfer}'s convention of treating the always policy
as the full-escalation endpoint; it is an analytical mixture, not a
separately run randomized arm. Non-commits and a sensitivity check
using the actual always-arm rate are reported in
Appendix~\ref{app:current-protocol}.
Timing is reconstructed from chunk offsets and logged stage durations
as time to first audio: the relay path ends when its synthesized
waveform is ready (expert wait plus blocking TTS), excluding the
waveform's playback duration; the local path ends at text generation.
Figure~\ref{fig:dualview} reports mean and P50 of this
diagnostic; Table~\ref{tab:latency} retains internal P95/P99.
These summaries do not measure client-audible completion
(Appendix~\ref{app:current-protocol}).


\paragraph{Main-table replay and open-ended scores.}
The NVDA lower block re-mixes local and cached expert outcomes after
pass-3 replay through the native frame protocol
(Appendix~\ref{app:nvda}). External columns retain cached local
outcomes and the prior 2,481-row fit. The Internal column judges the
same pass-3 answers on all 240 queries and uses a 2,258-row refit that
excludes these test IDs, with cached pass-2 training labels. All
policies leave the 17 queries without an answer-onset read local,
so 15/30/50\% scoreable-query budgets realize 14/28/47\% overall.
This Internal result is post-selection because probe architecture
selection included this split.

AlpacaEval uses a separate 1--5 scale: MiniCPM changes from 3.68
locally to 4.01 at 44\% escalation (rounded-anchor matched random
3.91; always-escalate 4.20, compared with 4.95 for returned expert
text). The pass-3 NVDA ranking changes from 3.55 to 4.51 versus
4.23 for random. These open-ended scores do not enter the QA macro
average.
